\section{Default theories and autoepistemic expansions: proofs}
\label{sec:annex-p18-default-autoepistemic}

\begin{annexitem}{annex:p18.default.finite.support}{thm:p18-default-finite-support}{Every closure and finite-support direction}
\begin{proof}
The stages \(C_n(B)\) increase because each stage is
contained in its classical closure. Their union contains
\(W\). A classical consequence of that union uses
finitely many premises; they all lie in a common
finite stage, so their consequence lies in the next.
A licensed default whose prerequisite lies in the
union has that prerequisite in a finite stage,
and contributes its conclusion at the next stage.
Thus the union has all the stated closure properties.
Any other set with those properties contains every
stage by induction, proving leastness. Increasing
\(B\) preserves each licensed default and hence
every stage, proving monotonicity.

Every finite support derivation lies in \(C(B)\):
background lines do, classical steps use its
deductive closure, and default steps use its
licensed closure. Conversely, induct on the
stage containing a given formula. A background
formula has empty support. A classical
consequence uses finitely many earlier formulas;
combine their finite derivations and union
their finite supports. A default step combines
the prerequisite's finite derivation with that
default and adds its finite justification set.
This gives a finite support in \(B\).
The argument handles infinite background sets
without assuming all their members occur in
one finite derivation.
\end{proof}
\end{annexitem}

\begin{annexitem}{annex:p18.default.extension.properties}{thm:p18-default-extension-properties}{Generating defaults and the exact consistency condition}
\begin{proof}
Let \(E=\Gamma_{\mathcal T}(E)\) and let \(G\)
be the displayed generating defaults. Every
new default conclusion at a construction
stage comes from \(G\), since its prerequisite
is already in \(E\) and its justifications
pass the final \(E\) test. Induction on
stages therefore places \(E\) inside the
displayed classical theory. Conversely every
generating prerequisite belongs to some
finite stage of \(E\), so its conclusion
appears at the next. All generating
conclusions and \(W\) therefore lie in
the closed \(E\), proving the reverse
inclusion.

If closed \(E\subseteq F\), consistency of
\(F\cup\{j\}\) implies consistency of
\(E\cup\{j\}\). Thus \(B(F)\subseteq B(E)\);
monotonicity of licensed closure proves
\(\Gamma(F)\subseteq\Gamma(E)\).
If both are extensions, \(E\subseteq F\)
therefore also implies \(F\subseteq E\).

The only inconsistent closed theory is
\(\mathcal L_0\). No individual justification
is consistent with it, so
\(\Gamma(\mathcal L_0)=C(\varnothing)\).
An inconsistent extension consequently
exists exactly when \(C(\varnothing)\)
is inconsistent. In that case every
candidate's licensed closure contains
the inconsistent \(C(\varnothing)\)
and is all of \(\mathcal L_0\);
there is only that one extension.
When each default has a nonempty
justification list, none is licensed
by the empty set, so
\(C(\varnothing)=\operatorname{Th}(W)\).
The no-extension example below proves
the last asserted limitation.
\end{proof}
\end{annexitem}

\begin{annexitem}{annex:p18.default.admissible.consistency}{thm:p18-default-admissible-consistency}{Admissible supports and complete flattening}
\begin{proof}
Write \(H=H(S)\). If \(C(\varnothing)\)
is inconsistent, so is its superset
\(C(H)\). Conversely suppose
\(C(\varnothing)\) is consistent but
\(\bot\in C(H)\). The finite-support
lemma supplies a finite \(K\subseteq H\)
supporting \(\bot\). This \(K\) is
nonempty, since empty support would
contradict consistency of \(C(\varnothing)\).
Choose \(j\in K\) and an argument
\((K_0,\varphi_0)\in S\) with \(j\in K_0\).
Explosion from the finite contradiction
gives the argument \((K,\bar j)\).
It attacks \((K_0,\varphi_0)\).
Admissibility supplies \((K_1,\psi)\in S\)
attacking \((K,\bar j)\), so
\(\bar\psi\in K\subseteq H\).
Some argument of \(S\) has
\(\bar\psi\) in its support and
is therefore attacked by that
same \((K_1,\psi)\). This violates
conflict-freeness of \(S\).
The consistency equivalence follows.
Every member of \(\operatorname{flat}(S)\)
belongs to \(C(H)\), proving the
flattening consistency consequence.

If \(S\) is complete, it contains
all empty-support arguments, since
they have no attackers and are
defended by every set. Hence its
flattening contains \(W\).
Suppose \(\varphi\) follows classically
from finitely many flattened
\(\varphi_1,\ldots,\varphi_m\), represented
by \((K_i,\varphi_i)\in S\).
Combining their finite derivations
gives an argument
\((K_1\cup\cdots\cup K_m,\varphi)\).
Every attacker of this combined
support attacks some original
\((K_i,\varphi_i)\), and is therefore
attacked by \(S\). The combined
argument is defended, so completeness
puts it in \(S\). The case of
no premises uses empty support.
Thus the flattening is classically
closed.
\end{proof}
\end{annexitem}

\begin{annexitem}{annex:p18.default.flattening}{thm:p18-default-flattening}{The full arg/flat identities}
\begin{proof}
The arguments in \(\operatorname{arg}(E)\)
are exactly finite supports contained
in \(B(E)\). Their conclusions are
therefore exactly \(C(B(E))\), by
the finite-support lemma. This
is \(\Gamma_{\mathcal T}(E)\).

A stable argument extension \(S\)
is complete by the general extension
relations, so \(E=\operatorname{flat}(S)\)
is closed by the preceding lemma.
An argument \((K,\varphi)\) is
not attacked by \(S\) exactly
when \(\bar j\notin E\) for
every \(j\in K\): its attacker
must conclude precisely \(\bar j\).
For closed \(E\), this is
equivalent to consistency of
each \(E\cup\{j\}\), and hence
to membership in \(\operatorname{arg}(E)\).
Stability is the fixed-point
equation for the no-attack map,
so \(S=\operatorname{arg}(E)\).
For inconsistent \(E\), every
nonempty support is attacked
and only the empty supports
remain; the same equivalence
still holds.
\end{proof}
\end{annexitem}

\begin{annexitem}{annex:p18.default.stable.bijection}{thm:p18-default-stable-bijection}{Both directions and the inverse maps}
\begin{proof}
If \(E\) is a Reiter extension,
the flattening identity gives
\(\operatorname{flat}(\operatorname{arg}(E))=E\).
The no-attack calculation in
the preceding proof, now applied
to this argument set, says an
argument is not attacked by it
exactly when it belongs to
\(\operatorname{arg}(E)\).
Thus this set is stable.
Conversely, for stable \(S\),
the same identities give
\[
 \operatorname{flat}(S)
 =\operatorname{flat}(\operatorname{arg}(\operatorname{flat}(S)))
 =\Gamma_{\mathcal T}(\operatorname{flat}(S)).
\]
Its flattening is therefore
a Reiter extension. The two
already proved compositions
are identity maps on the
respective families.
Stable extensions are preferred,
and the arbitrary-framework
maximality theorem supplies
preferred extensions whether
or not the stable family
is nonempty.
\end{proof}
\end{annexitem}

\begin{annexitem}{annex:p18.default.examples}{thm:p18-default-examples}{The complete source example and separating cases}
\begin{proof}
For \(\mathcal T_{\rm paradox}\),
the default is licensed by a
candidate exactly when \(p\)
is absent from that candidate.
Its prerequisite is always
available. Thus \(\Gamma(E)\)
is \(\operatorname{Th}(\{p,q\})\)
when \(p\notin E\), and
\(\operatorname{Th}(\{q\})\)
when \(p\in E\). Neither
can equal its corresponding
candidate. There is no
Reiter extension, and the
bijection gives no stable
argument extension.

The only possible supports
are empty and \(\{\neg p\}\).
Their licensed consequences
are respectively
\(\operatorname{Th}(\{q\})\)
and \(\operatorname{Th}(\{p,q\})\).
The argument
\(b=(\{\neg p\},p)\)
attacks every nonempty-support
argument and itself.
No admissible set contains
an argument concluding \(p\),
since any such argument
has that support and
self-attacks. Therefore no
admissible set can defend
any nonempty-support argument
against \(b\): attacking
\(b\) requires a conclusion
\(p\). All admissible
sets are subsets of the
empty-support set \(S_0\).
Every member of \(S_0\)
is unattacked, and none
concludes \(p\); hence
\(F(S_0)=S_0\). It is
the unique preferred and
complete extension, and
is grounded. Its flattening
is exactly \(\operatorname{Th}(\{q\})\).

For \(\mathcal T_{\rm separate}\),
the only possible generating
conclusion is \(r\). The
candidate \(\operatorname{Th}(\{r\})\)
is individually consistent
with \(p\) and with \(\neg p\),
so is fixed; the candidate
\(\operatorname{Th}(\varnothing)\)
activates the default and
is not fixed. The generating
default theorem shows there
are no other candidates.
No argument concludes \(p\)
or \(\neg p\): licensed
closures are either tautologies
or consequences of \(r\).
Those are the only possible
conclusions which could attack
a justification, so the
entire graph is unattacked.
Its whole argument set
is its unique stable
extension and flattens
to \(\operatorname{Th}(\{r\})\).

For \(\mathcal T_{\rm self}\),
even licensing its justification
cannot derive the prerequisite
\(p\) from the empty background.
Every licensed closure is
therefore \(\operatorname{Th}(\varnothing)\),
which is the unique extension.
In the inconsistent-background
case explosion holds already
in the unconditional closure.
In the justification-free case
the unconditional rule derives
\(\bot\), with the same effect.
The consistency theorem proves
the asserted unique inconsistent
extension in both cases.
\end{proof}
\end{annexitem}

\begin{annexitem}{annex:p18.default.finite.method}{thm:p18-default-finite-method}{Finite propositional enumeration}
\begin{proof}
There are finitely many subsets
of \(D\), and every candidate
has a finite propositional basis.
Consequence and consistency of
such bases are decidable by
their finite truth assignments.
Once the candidate's consistency
tests are fixed, the licensed
loop adds only conclusions
from \(D\). It therefore
has at most \(|D|\) strict
enlargements. At termination
its closed theory is exactly
the least licensed closure,
by the closure lemma.
Two finite bases have the
same consequences exactly
when each entails every
member of the other, which
the same finite tests decide.

A returned candidate is
therefore exactly a fixed
point of \(\Gamma\).
Conversely every extension
is generated by a subset
of the defaults, by the
generating-default theorem,
so its candidate occurs
and passes. Duplicate
bases are removed by the
same equivalence test.
These arguments use finite
propositional inputs only;
they do not convert general
first-order consistency into
a decision procedure.
\end{proof}
\end{annexitem}

\begin{annexitem}{annex:p18.ael.kernel}{thm:p18-ael-kernel}{The full modal-depth completion and uniqueness proof}
\begin{proof}
First suppose \(E\) is consistent.
The recursion is well-defined:
in evaluating an outer \(L\psi\),
membership of \(\psi\) has strictly
smaller modal depth. It assigns
a fixed truth value to every
closed modal atom. A classical
model of \(E\) can be expanded
by assigning all modal atoms
those fixed values. For every
sentence \(\theta\), induction
on formula construction says
its value in that expansion
is the value of \(\theta^E\)
in the original model.
The Boolean cases are their
truth tables. First-order
quantifiers commute with the
replacement because each
replaced modal atom is closed
and independent of the bound
variable. This is exactly
where the absence of
quantifying-in is used.

Consequently the sentences
true in all these expanded
models are precisely
\(\operatorname{St}(E)\).
This set is classically
closed and has objective
part \(E\). If \(\psi\)
belongs to it, its assigned
\(L\psi\) value is true;
otherwise \(\neg L\psi\)
is assigned true. Thus it
contains its full
\({\cal B}(\operatorname{St}(E))\).
The models of that assignment
together with \(E\) are
exactly the expansions just
described, proving the
displayed consequence equation.

For uniqueness, let \(T\)
be closed, contain
\({\cal B}(T)\), and have
objective part \(E\).
For each sentence \(\theta\),
finitely many members of
\({\cal B}(T)\) fix the
outer modal atoms occurring
in it. Replacing them by
their assigned constants
gives an objective sentence
\(\theta^T\) with
\(\theta\leftrightarrow\theta^T\)
a consequence of those
finitely many assignments.
Hence \(\theta\in T\)
exactly when \(\theta^T\in E\).
Inducting on modal depth
now shows these assignments
agree with the construction
from \(E\), and hence
\(T=\operatorname{St}(E)\).
This proves uniqueness for
all nested modalities, not
only depth-one formulas.

If \(E\) is inconsistent,
its classical closure is
all of \(\mathcal L_0\).
The recursion consequently
includes every sentence,
and any closed theory
with that objective part
is all of \(\mathcal L_L\).
All asserted equations
and uniqueness then follow
by explosion. Finally a
stable expansion is closed
and contains its introspective
assignment by its defining
equation, so the uniqueness
result applies to it.
\end{proof}
\end{annexitem}

\begin{annexitem}{annex:p18.default.ael.weak}{thm:p18-default-ael-weak}{The exact weak-extension correspondence}
\begin{proof}
Let \(T\) be an expansion and
\(E=T\cap\mathcal L_0\).
The kernel theorem gives
\(T=\operatorname{St}(E)\).
Under its fixed modal
assignment, \(L\alpha_d\)
is true exactly when
\(\alpha_d\in E\), and
\(\neg L\bar\beta_{d,i}\)
is true exactly when
\(\bar\beta_{d,i}\notin E\).
Thus each transformed
default is equivalent,
under that assignment,
either to its conclusion
\(\gamma_d\) if all its
tests pass, or to
\(\top\) otherwise.
This equivalence holds
also for an empty
justification conjunction.
It follows that
\[
 \operatorname{Cn}_L(\tau(\mathcal T)\cup{\cal B}(T))
 =\operatorname{Cn}_L(\Delta_{\mathcal T}(E)\cup{\cal B}(T)).
\]
Taking objective parts
gives \(E=\Delta_{\mathcal T}(E)\).
The reduction remains
valid for an inconsistent
kernel: all positive modal
atoms are true, so only
the zero-justification
conclusions pass.

Conversely, if closed
\(E=\Delta_{\mathcal T}(E)\),
put \(T=\operatorname{St}(E)\).
The same modal reduction
and the kernel consequence
equation give
\[
 \operatorname{Cn}_L(\tau(\mathcal T)\cup{\cal B}(T))
 =\operatorname{Cn}_L(E\cup{\cal B}(T))=T.
\]
Thus \(T\) is an expansion.
The kernel and completion
maps are inverse.
The generating-default
identity already proved
for every Reiter extension
is exactly \(E=\Delta(E)\).
This proves the general
forward implication.
The explicit self-support
example below refutes
its unrestricted converse.
\end{proof}
\end{annexitem}

\begin{annexitem}{annex:p18.default.ael.prerequisite.free}{thm:p18-default-ael-prerequisite-free}{Both directions for prerequisite-free theories}
\begin{proof}
Fix a candidate \(E\).
All defaults licensed by
\(B(E)\) have prerequisite
\(\top\), which belongs
to the classical closure
of every background.
Their conclusions therefore
all enter the licensed
construction, regardless
of whether any other
default has fired.
Its least closed result
is exactly the classical
closure of \(W\) and
these conclusions.
This is \(\Delta(E)\),
since every closed
candidate contains \(\top\).
The two fixed-point
conditions coincide.

Combine the weak-extension
bijection with Dung's
arg/flat bijection to
obtain the three inverse
correspondences. Every
introspective closed theory
contains \(L\top\), so
its omission from the
transformed antecedent
changes no expansion.
All derivations used
finite proof witnesses
for individual formulas;
no finiteness of the
input theory was assumed.
The equality of operators
also applies to the
inconsistent candidate,
so that case is included.
\end{proof}
\end{annexitem}

\begin{annexitem}{annex:p18.default.ael.boundary}{thm:p18-default-ael-boundary}{The exact self-support counterexample}
\begin{proof}
For the one-default theory,
\(\Delta(E)\) is either
\(\operatorname{Th}(\{p\})\)
if \(p\in E\) and
\(\neg p\notin E\), or
\(\operatorname{Th}(\varnothing)\)
otherwise. These are
therefore the only
possible weak fixed
points. Both are fixed:
the first activates the
default, and the second
does not satisfy its
prerequisite. Neither
is inconsistent; the
inconsistent candidate
would disable the
justification and is
not fixed.

The weak-extension theorem
therefore gives exactly
the two asserted ordinary
stable expansions. The
default-example proof
showed that Reiter's
construction never derives
\(p\) from its empty
background, even when
its justification is
licensed. Thus only
the tautological kernel
is a Reiter extension.
This is a counterexample
to the general bijection
for this direct transform.
It proves no impossibility
claim about other, more
involved translations.
\end{proof}
\end{annexitem}
